% ============================================================
%  Section: Experimental Results
%  Writing order: SECOND (write after Methods)
% ============================================================
\section{Experimental Results}
\label{sec:results}
% TODO: Write results.
% Recommended subsection structure:
%
%   \subsection{Datasets}
%     - Table: dataset name, N subjects, channels, sampling rate, label balance
%     - MODMA: N~53, note small-N implication for deep networks
%     - Mumtaz / other datasets used
%
%   \subsection{Within-Dataset Performance (Nr\,DE, K=3)}
%     - Table: per-dataset AUC mean +/- std across LOSO folds
%     - Claims.csv entries: result claim for K=3 nrde AUC
%     - Ablation: K=1,2,3,5,10 AUC curve (if available)
%
%   \subsection{Cross-Dataset Performance (SpecSPD-DANN)}
%     - Table: source -> target AUC, mean +/- std, N folds
%     - Claims.csv entries: result claim for SpecSPD-DANN cross-dataset
%     - Seed stability analysis (3-5 seeds; report std across seeds)
%
%   \subsection{Ablation: Deep Heads vs.\ Riemannian Features}
%     - Motivation: deep networks overfit at N~53 (see claims.csv BG-003)
%     - Show AUC degradation with MLP/CNN heads on MODMA
%
%   \subsection{Comparison with Prior Work}
%     - Table: method, dataset, evaluation protocol, AUC/acc reported
%     - IMPORTANT: only compare methods using equivalent protocols
%       (no apples-to-oranges with within-subject results) — CLAUDE.md Rule 7
%
% Every numeric claim in this section MUST have a claims.csv row
% with evidence_file pointing to results/*.json or results/*.log.
